\section{Density Zeros and the Linear Equivalence Group}
\label{sec:density-zeros-and-the-linear-equivalence-group}

The assumption \(f>0\) Lebesgue-almost everywhere has a different role
from the moment and Sobolev assumptions.
It makes finite changes of coordinates equivalent and supplies the
translation quasi-invariance used in the zero--one argument.
Removing it changes the classification itself.
In this section the hypotheses are stated separately for each result;
\eqref{eq:main-assumptions} is not a standing assumption.
The rotational Fourier defect still detects non-Gaussian mixing, but a finite
Hellinger cost need not represent a change within the same measure class.

\subsection{Failure of dichotomy without density positivity}
\label{sec:restricted-counterexample}

Let \(f\) be the variance-one rescaling of
\(C(1-x^2)^2\mathbf 1_{\{|x|<1\}}\), and write its support as
\([-a,a]\).
Set \(\rho(dx)=f(x)\,dx\).
Its zero-extended square root belongs to \(W^{1,2}(\R)\), and
\(x(\sqrt f)'\in L^2(\R)\).
For \(0<c<1\), set
\[
 T=c\oplus I,\qquad
 f_c(y)=c^{-1}f(y/c),\qquad \mu=\rho^{\otimes\N}.
\]
Then \(T-I\) has rank one and
\[
 \nu_T=(f_c(y)\,dy)\otimes\rho^{\otimes\{2,3,\ldots\}}\ll\mu,
 \qquad \mu\not\ll\nu_T.
\]
Indeed, \(f_c\) is positive on \((-ca,ca)\) and zero on the rest
of \((-a,a)\).
Nevertheless,
\[
 \Aff(\mu,\nu_T)=\int\sqrt{f(y)f_c(y)}\,dy>0,
 \qquad \sup_n E_n(T)<\infty.
\]
Thus neither Hilbert--Schmidt proximity nor bounded projected energy
implies equivalence without an additional condition on null sets.
The failure occurs with finite location and scale Fisher information.

\subsection{What survives of the Fourier and Fisher argument}
\label{sec:restricted-local}

The logarithm-free Fourier characterization in
\cref{lem:fourier-ode}, and its one-rotation Hellinger lower bound,
already require no density.
For the stronger matrix estimates, the appropriate regularity condition is
\begin{equation}
 \begin{gathered}
 f\geq0,\quad \int f=1,\quad
 g=\sqrt f\in W^{1,2}(\R),\quad xg'\in L^2(\R),\\
 \E X=0,\quad \E X^2=1,\quad f\text{ non-Gaussian}.
 \end{gathered}
 \label{eq:restricted-sobolev}
\end{equation}
Here \(g\) is defined on the whole real line, including by zero outside
the support.
Regularity only in the interior of the support is insufficient.
For example, the uniform density has zero interior derivative but its
square root has jumps at the boundary.

\begin{proposition}[Local Fourier coercivity with zeros of the density]
\label{prop:restricted-local}
Under \eqref{eq:restricted-sobolev}, the matrix Fisher decomposition,
finite-dimensional Hellinger path bound, and dimension-free local
Frobenius coercivity in \cref{sec:finite-dimensional-hellinger-coercivity} remain valid.
Write \(\mu=\rho^{\otimes\N},\quad\rho(dx)=f(x)\,dx\).
There are constants \(c,C,\delta>0\), with \(\delta<1\), such that
\begin{equation}
 c\|K\|_{\mathcal S_2}\leq h(\mu,\nu_{I+K})
       \leq C\|K\|_{\mathcal S_2}
 \qquad(\|K\|_{\mathcal S_2}\leq\delta).
 \label{eq:restricted-hs-local}
\end{equation}
Moreover, for every \(T\in\GL(H)\), the following necessary
condition survives:
\begin{equation}
 \nu_T\ll\mu
 \quad\Longrightarrow\quad
 T-Q\in\mathcal S_2
 \quad\text{for some }Q\in\mathcal P.
 \label{eq:restricted-hs-necessary}
\end{equation}
For orthogonal transformations one may omit \(xg'\in L^2\),
retaining non-Gaussianity, \(g\in W^{1,2}(\R)\), and the centered
variance-one normalization.
These statements do not assert absolute continuity of
\(\nu_{I+K}\).
\end{proposition}
\begin{proof}
A Sobolev function has weak derivative zero almost everywhere on its
zero set.
Define \(s=-2g'/g\) on \(\{g>0\}\) and \(s=0\) elsewhere.
Then \(sg=-2g'\) almost everywhere on \(\R\).
The cutoff integration-by-parts identities, including
\(\E s(X)=0\), \(\E[Xs(X)]=1\), and the formulas for \(J,\kappa\), are
therefore unchanged.
Equality \(J=1\) still gives the global equation \(g'=-xg/2\);
it cannot produce a Gaussian truncated at a boundary.
Similarly, \(\kappa=0\) would give \(2xg'+g=0\) globally and no
nonzero square-integrable probability half-density.
The finite real Fourier sums with dual moments and half-density generator estimates use
these identities and \(L^2\) regularity, so their proofs apply.

The infinite-dimensional limiting step needs a different interpretation.
For \(K\in\mathcal S_2\) with \(I+K\) invertible, put
\(T_n=I+P_nKP_n\).
The finite-dimensional path bound and invariance of Hellinger distance
give, for sufficiently large \(m,n\),
\[
 h(\nu_{T_n},\nu_{T_m})
 \leq C\|T_n^{-1}\|_{\mathrm{op}}
                 \|T_n-T_m\|_{\mathcal S_2}.
\]
Choose a single probability measure dominating the countable family
\(\mu,\nu_{T_1},\nu_{T_2},\ldots\).
Their square-root densities are Cauchy in its \(L^2\) space.
The nonnegative limit has norm one and defines a probability measure;
convergence of finite output row sums identifies this measure with
\(\nu_{I+K}\).
This proves Hellinger convergence, not a density relative to \(\mu\).
Passing the finite local bounds to this limit proves
\eqref{eq:restricted-hs-local}.
For orthogonal operators, use the finite-coordinate polar
approximants and skew-generator bound from
\cref{prop:orthogonal-sufficiency} in the same way.

Finally, the proof of \cref{lem:hilbert-schmidt-necessity} uses absolute
continuity only through
\cref{lem:tail-symmetry-uniform,lem:compact-covariance-rigidity},
which assume \(\nu_T\ll\mu\).
Its other inputs are the moment-only characteristic-function rigidity,
the measurable inverse action, and the local estimate for the
finite-rank differences \(PT^{-1}ST-I\).
The estimate just established supplies this last input without any
quasi-invariance assumption.
The same tail matching and permutation argument therefore proves
\eqref{eq:restricted-hs-necessary}.
\end{proof}

For standardized Gamma shape \(\alpha\), the zero-extended square
root is in \(W^{1,2}(\R)\) exactly when \(\alpha>2\).
For an affine rescaling of a Beta\((\alpha,\beta)\) density the
corresponding condition is \(\alpha,\beta>2\).
These assertions follow by squaring the derivative of the boundary
power \(t^{(\alpha-1)/2}\).
At shape one the square root instead has a jump, so its vanishing
interior power derivative is not an exception.
The support classifications below require none of this regularity
and include every positive Gamma or Beta shape.

\subsection{The linear equivalence group for arbitrary zero sets of the density}
\label{sec:linear-equivalence-group}

The Hilbert--Schmidt condition and the set where \(f\) is positive together
give an exact operator criterion, including for densities with gaps and
with both tails unbounded.  The required condition on the rows concerns the
actual zero set of the density, rather than its closed support.

\begin{theorem}[The linear equivalence group with density zeros]
\label{thm:linear-equivalence-group}
Assume \eqref{eq:restricted-sobolev}, let $\rho(dx)=f(x)\,dx$, and put
$\mathcal U=\{x\in\R:f(x)>0\}$.  For every $T\in\GL(H)$,
\begin{equation}
 \nu_T\ll\mu
 \quad\Longleftrightarrow\quad
 \begin{cases}
 T-Q\in\mathcal S_2(H)\text{ for some }Q\in\mathcal P,\\
 \displaystyle
 \mu\!\left\{x:\sum_jT_{ij}x_j\notin\mathcal U\right\}=0
       \text{ for every }i.
 \end{cases}
 \label{eq:absolute-continuity-operator-criterion}
\end{equation}
Consequently the linear equivalence group of $\mu$ is
\begin{equation}
 \begin{aligned}
 \mathcal G_\rho
 &:=\{T\in\GL(H):\nu_T\sim\mu\}\\
 &=\left\{T\in\GL(H):
 \begin{array}{l}
 T-Q\in\mathcal S_2(H)\text{ for some }Q\in\mathcal P,\\
 \displaystyle\mu\!\left\{x:\sum_jT_{ij}x_j\notin\mathcal U\right\}=0,\\
 \displaystyle\mu\!\left\{x:\sum_j(T^{-1})_{ij}x_j\notin\mathcal U\right\}=0
       \quad\text{for every }i
 \end{array}\right\}.
 \end{aligned}
 \label{eq:linear-equivalence-group}
\end{equation}
This set is a subgroup of $\GL(H)$.  For orthogonal members of
$\mathcal G_\rho$, the same criterion holds with $T^{-1}=T^*$
and without the scale condition $xg'\in L^2$.

More precisely, whenever $T-Q\in\mathcal S_2(H)$ for some
$Q\in\mathcal P$, the Lebesgue decomposition of $\nu_T$ relative to
$\mu$ is its restriction to $\mathcal U^{\N}$ plus its restriction
to the complement: the former is absolutely continuous and the latter
is singular.
\end{theorem}
\begin{proof}
First suppose $T-Q\in\mathcal S_2$ and put $K=TQ^{-1}-I$.
The exact invariance $\nu_Q=\mu$ gives $\nu_T=\nu_{I+K}$.
For the eventually invertible finite-coordinate approximants
$T_n=I+P_nKP_n$, the proof of \cref{prop:restricted-local} gives
\[
 h(\nu_{T_n},\nu_T)\longrightarrow0.
\]
In particular their probabilities converge uniformly over measurable
sets, by the Hellinger bound on total variation.

For each such $n$, the law $\nu_{T_n}$ has a Lebesgue density in the
first $n$ coordinates and retains the independent product law on the
remaining coordinates.  On $\mathcal U^n$, the density of $\mu_n$ is
strictly positive.  Fubini's theorem therefore gives
\[
 \left.\nu_{T_n}\right|_{\mathcal U^{\N}}\ll\mu.
\]
If $B$ is any $\mu$-null Borel set, then
$\nu_{T_n}(B\cap\mathcal U^{\N})=0$ for all sufficiently large $n$.
Total variation convergence implies
$\nu_T(B\cap\mathcal U^{\N})=0$.
Thus $\left.\nu_T\right|_{\mathcal U^{\N}}\ll\mu$.
Since $\mu(\mathcal U^{\N})=1$, the restriction to its complement is
singular.  This proves the asserted decomposition and
\[
 \nu_T\ll\mu
 \quad\Longleftrightarrow\quad
 \nu_T(\mathcal U^{\N})=1
 \quad\Longleftrightarrow\quad
 \mu\!\left\{x:\sum_jT_{ij}x_j\notin\mathcal U\right\}=0
       \text{ for every }i.
\]
The necessary Hilbert--Schmidt condition in
\eqref{eq:restricted-hs-necessary} now proves
\eqref{eq:absolute-continuity-operator-criterion} in both directions.

If $\nu_T\sim\mu$, the measurable inverse action gives
$\nu_{T^{-1}}\sim\mu$, so both row conditions in
\eqref{eq:linear-equivalence-group} are necessary.
Conversely, $T-Q\in\mathcal S_2$ implies
\[
 T^{-1}-Q^{-1}=-T^{-1}(T-Q)Q^{-1}\in\mathcal S_2.
\]
Apply \eqref{eq:absolute-continuity-operator-criterion} to $T$ and
$T^{-1}$ to obtain $\nu_T\ll\mu$ and $\nu_{T^{-1}}\ll\mu$.
The measurable inverse action then gives
\[
 \mu=(\Phi_T)_\#\nu_{T^{-1}}\ll(\Phi_T)_\#\mu=\nu_T,
\]
which proves equivalence.
For $S,T\in\mathcal G_\rho$, measurable composition yields
$\nu_{ST}=(\Phi_S)_\#\nu_T\sim\nu_S\sim\mu$; the inverse argument
already proved gives $T^{-1}\in\mathcal G_\rho$.
Thus the displayed operator class is a group.
For orthogonal transformations, repeat the restriction argument with
the finite-coordinate orthogonal approximants from
\cref{prop:orthogonal-sufficiency}; their Hellinger convergence needs
only location Fisher information.
\end{proof}

Every nonzero row sum has an absolutely continuous scalar law: condition
on all but one input with a nonzero coefficient.  Thus the row conditions
are unchanged if $f$ is modified on a Lebesgue-null set.
If $f>0$ almost everywhere, all these conditions are automatic and
\eqref{eq:linear-equivalence-group} reduces to
\cref{thm:measure-equivalence}.
For a density with gaps, they test whether each row of $T$ and of its
inverse puts any mass in those gaps.  No likelihood-ratio limit is
required in addition to these operator conditions.  The contraction in
\cref{sec:restricted-counterexample} satisfies the forward row condition
and fails the inverse row condition, giving strict absolute continuity.

\subsection{Complete linear classification for bounded and half-line laws}
\label{sec:restricted-support-rigidity}

Write \(\rho=\Law(X)\), \(\rho^-=\Law(-X)\), and
\(\mu=\rho^{\otimes\N}\).
Throughout this subsection \(X\) is centered with variance one.
The operator action is the row-series action of
\cref{lem:measurable-action}.

\begin{theorem}[Linear equivalence for bounded and half-line coordinate laws]
\label{thm:restricted-support-rigidity}
Let \(T\in\GL(H)\).
\begin{enumerate}
 \item Suppose \(a=\operatorname{ess\,inf}X<0<
 b=\operatorname{ess\,sup}X\) are finite.
 Then \(\nu_T\sim\mu\) forces \(T\) to be a signed permutation.
 The complete equivalence group is as follows:
 \[
 \begin{array}{c|l}
 \text{relation between }\rho^-\text{ and }\rho&\text{allowed signed permutations}\\ \hline
 \rho^-=\rho&\text{all signs}\\
 \rho^-\sim\rho,\ \rho^-\ne\rho&
                  \text{only finitely many negative signs}\\
 \rho^-\not\sim\rho&\text{no negative signs}.
 \end{array}
 \]
 In particular, \(a+b\ne0\) permits only ordinary permutations.
 The support need not fill an interval, and a density is not required.
 \item Suppose \(\operatorname{ess\,inf}X=a>-\infty\) and
 \(\operatorname{ess\,sup}X=+\infty\).
 Then
 \begin{equation}
  \nu_T\sim\mu
   \quad\Longleftrightarrow\quad
  T\text{ is an ordinary coordinate permutation}.
  \label{eq:restricted-halfline-classification}
 \end{equation}
 The same conclusion holds for a law bounded above and unbounded below.
\end{enumerate}
In the second case every equivalence is exact invariance.
Failure of any of these equivalence criteria does not, by itself,
imply singularity.
\end{theorem}
\begin{proof}
Assume \(\nu_T\sim\mu\), and set \(S=T^{-1}\).
The measurable inverse action implies
\((\Phi_S)_\#\nu_T=\mu\).
Pushing equivalence forward by \(\Phi_S\) therefore gives
\(\nu_S\sim\mu\).
For any row \(r\) of \(T\) or \(S\), the random series
\(Z=\sum_jr_jX_j\) converges in \(L^2\).
Centering and independence imply
\begin{equation}
 \E[Z\mid X_1,\ldots,X_m]=\sum_{j\leq m}r_jX_j.
 \label{eq:restricted-row-conditioning}
\end{equation}

In the bounded case \(Z\in[a,b]\) almost surely.
Its conditional expectation lies in the same interval.
Independence and the definition of essential endpoints show that
the essential width of the finite sum on the right of
\eqref{eq:restricted-row-conditioning} is
\((b-a)\sum_{j\leq m}|r_j|\).
Consequently every row of both \(T\) and \(S\) has
\(\ell^1\)-norm at most one.
This proves the required bound even when the original row has
infinitely many nonzero entries.

Let \(A_{ij}=|T_{ij}|\) and \(B_{ij}=|S_{ij}|\).
The products defining the entries of \(TS\) and \(ST\) are absolutely
convergent by the \(\ell^2\) row/column Cauchy--Schwarz inequality.
Thus \(AB\geq I\) and \(BA\geq I\) entrywise.
Tonelli's theorem and the row bounds give
\[
 \sum_j(AB)_{ij}
 =\sum_k A_{ik}\sum_jB_{kj}\leq1,
 \qquad
 \sum_j(BA)_{ij}\leq1.
\]
It follows that \(AB=BA=I\).
This reasoning concerns nonnegative matrix entries and does not
assume that entrywise absolute values define bounded operators
on \(\ell^2\).

For completeness, nonnegative matrices inverse to one another in
this sense are monomial.
For each \(i\), choose \(k\) with \(A_{ik}B_{ki}>0\).
The off-diagonal entries of \(AB=I\) force row \(k\) of \(B\)
to vanish outside column \(i\).
Then \(BA=I\) forces row \(i\) of \(A\) to vanish outside column \(k\).
The resulting correspondence is bijective.
Its weights and their reciprocals are at most one by the row bounds,
so all weights are one.
Hence \(T\) is a signed permutation.

For a negative row sign, scalar equivalence requires
\(\rho^-\sim\rho\), which in particular requires \(a=-b\).
The remaining classification is a product calculation.
Put \(r=\Aff(\rho,\rho^-)\).
When these laws are equivalent, \(r>0\), and \(r=1\) exactly when
they are equal.
Each negative sign contributes the factor \(r\), while positive
signs contribute one.
Kakutani's theorem~\cite{kakutani1948equivalence} gives precisely
the three alternatives stated above.

For the lower-bounded, upper-unbounded case, centering and
nondegeneracy give \(a<0\).
Now every finite sum in \eqref{eq:restricted-row-conditioning}
is bounded below by \(a\).
A negative coefficient would make its essential infimum
\(-\infty\), which is impossible.
Thus every \(r_j\geq0\), and its finite essential infimum gives
\[
 a\sum_{j\leq m}r_j\geq a,
 \qquad\text{hence}\qquad \sum_jr_j\leq1.
\]
Both \(T\) and \(S\) are therefore nonnegative with row sums at
most one.
The same inverse-matrix argument makes \(T\) an ordinary
permutation.
Such a permutation preserves the i.i.d. product, proving the
converse.
Reflection proves the upper-bounded case.
\end{proof}

For \(Y\sim\operatorname{Gamma}(\alpha,\theta)\), where \(\theta\)
is the scale, put
\[
 X=\frac{Y-\alpha\theta}{\sqrt{\alpha}\theta}.
\]
Its essential endpoints are \(-\sqrt{\alpha}\) and \(+\infty\).
Thus \eqref{eq:restricted-halfline-classification} gives the complete
linear equivalence classification for every \(\alpha,\theta>0\),
including the shape-two Fourier example in
\cref{sec:variance-and-affine-normalization}.
For compactly supported densities positive almost everywhere inside
\((-a,a)\), reflection equivalence is automatic.
An asymmetric such density permits finitely many sign changes,
although they do not preserve the measure exactly.
A standardized asymmetric Beta law has unequal endpoint magnitudes,
so it permits only permutations.

\subsection{Exact block criteria without equivalence of the factors}
\label{sec:restricted-blocks}

For block products, the following criterion separates conditions for absolute
continuity within each pair of factors from accumulation of Hellinger loss.

\begin{theorem}[Block products without equivalence of the factors]
\label{thm:block-kakutani}
Let \(\mu=\bigotimes_{k\geq1}\lambda_k\) and
\(\nu=\bigotimes_{k\geq1}\eta_k\) be probability products over
finite-dimensional coordinate blocks.
Set
\[
 r_k=\Aff(\lambda_k,\eta_k),\qquad
 L=\sum_k-\log r_k\in[0,\infty],
 \qquad -\log0=+\infty.
\]
Then
\begin{equation}
 \Aff(\mu,\nu)=\prod_kr_k,\qquad
 \mu\perp\nu\quad\Longleftrightarrow\quad L=\infty.
 \label{eq:restricted-block-affinity}
\end{equation}
If \(L<\infty\), then
\begin{equation}
 \begin{aligned}
  \nu\ll\mu&\quad\Longleftrightarrow\quad
                         \eta_k\ll\lambda_k\text{ for every }k,\\
  \mu\ll\nu&\quad\Longleftrightarrow\quad
                         \lambda_k\ll\eta_k\text{ for every }k.
 \end{aligned}
 \label{eq:restricted-block-ac}
\end{equation}
Thus positive affinity allows equivalence, either direction of strict
absolute continuity, or neither direction.
For a linear block \(A_k\in\GL(d_k)\) and
\(\lambda_k=\rho^{\otimes d_k}\), where \(\rho(dx)=f(x)\,dx\),
the comparison density is
\[
 |\det A_k|^{-1}\prod_{i=1}^{d_k}
                         f((A_k^{-1}y)_i).
\]
The local absolute-continuity conditions concern its actual zero
set, not just the closure of its support.
\end{theorem}
\begin{proof}
Affinity tensorizes on each finite block prefix.
Convergence of projected affinities to global affinity follows by
conditional expectation under the common dominating measure
\((\mu+\nu)/2\); the proof is also given below.
This gives \eqref{eq:restricted-block-affinity}.

Necessity of each local absolute-continuity condition follows by
projection.
For sufficiency, suppose \(\eta_k\ll\lambda_k\) for all \(k\) and
\(\prod_kr_k>0\).
With \(\ell_n=\prod_{k\leq n}d\eta_k/d\lambda_k\), independence gives
\[
 \|\sqrt{\ell_m}-\sqrt{\ell_n}\|_{L^2(\mu)}^2
       =2\left(1-\prod_{n<k\leq m}r_k\right)\longrightarrow0.
\]
The square roots have norm one, so their limit squared is a
probability density.
Cylinder tests identify its measure with \(\nu\).
Interchanging \(\mu,\nu\) proves the reverse assertion.
This is the product argument behind Kakutani's criterion; the
mutual absolute-continuity hypothesis on each factor cannot be
discarded~\cite{kakutani1948equivalence,shepp1965translate}.
\end{proof}

\begin{corollary}[Masses of the absolutely continuous parts]
\label{cor:block-absolutely-continuous-masses}
Take \(m_k=(\lambda_k+\eta_k)/2\), with densities \(p_k,q_k\), and put
\[
 a_k=\lambda_k(\{q_k>0\}),\qquad
 b_k=\eta_k(\{p_k>0\}).
\]
If \(L<\infty\), denote the absolutely continuous parts of \(\nu\)
relative to \(\mu\), and of \(\mu\) relative to \(\nu\), by
\(\nu_a\) and \(\mu_a\), respectively.
Then
\begin{equation}
 \nu_a(\R^\N)=\prod_kb_k,\qquad
 \mu_a(\R^\N)=\prod_ka_k.
 \label{eq:restricted-common-masses}
\end{equation}
\end{corollary}
\begin{proof}
Since \(r_k^2\leq a_kb_k\) and \(L<\infty\), both displayed products are positive.
On \(C=\prod_k\{p_kq_k>0\}\), the two conditional product laws have
equivalent factors and affinities \(r_k/\sqrt{a_kb_k}\) with positive
product; hence they are equivalent.
Outside \(C\), each measure is carried by a set null for the other.
Their restrictions to \(C\) are therefore exactly their absolutely
continuous parts.
\end{proof}

\paragraph{Gamma scale changes and the effect of centering.}
For the uncentered shape-\(\alpha\), scale-one Gamma density, positive
scaling by \(c\) preserves the scalar measure class and gives
\begin{equation}
  \int_0^\infty\sqrt{f(y)c^{-1}f(y/c)}\,dy
 =\left(\frac{2\sqrt c}{1+c}\right)^\alpha
 =\cosh\!\left(\frac{\log c}{2}\right)^{-\alpha}.
 \label{eq:restricted-gamma-scale}
\end{equation}
Integration of \(y^{\alpha-1}e^{-(1+c^{-1})y/2}\) proves the formula.
Since \(\log\cosh(u/2)\sim u^2/8\) at zero, it follows that
\begin{equation}
 \bigotimes_k\Law(c_kY)\sim\bigotimes_k\Law(Y)
 \quad\Longleftrightarrow\quad
 \sum_k(\log c_k)^2<\infty,
 \qquad c_k>0.
 \label{eq:restricted-gamma-scale-product}
\end{equation}
The condition is equivalently \(\sum_k(c_k-1)^2<\infty\).
This does not contradict \cref{thm:restricted-support-rigidity}:
for \(X=(Y-\alpha)/\sqrt\alpha\), the map \(Y\mapsto cY\) becomes
\[
 X\longmapsto cX+\sqrt\alpha(c-1),
\]
an affine map fixing the lower endpoint, not the homogeneous
linear map \(X\mapsto cX\).

\begin{corollary}[Affine Gamma classification]
\label{cor:restricted-gamma-affine}
Let \(\mu\) be the standardized shape-\(\alpha\) Gamma product.
For \(T\in\GL(H)\) and \(b\in\ell^2\), the law of
\(\Phi_T(x)+b\) under \(\mu\) is equivalent to \(\mu\) if and only if
there are a coordinate
permutation \(\pi\) and numbers \(c_i>0\) such that
\[
 (Tx)_i=c_ix_{\pi(i)},\qquad
 b_i=\sqrt\alpha(c_i-1),\qquad
 \sum_i(\log c_i)^2<\infty.
\]
\end{corollary}
\begin{proof}
The inverse affine map has linear part \(T^{-1}\) and shift
\(-T^{-1}b\in\ell^2\).
To justify this for the row-series versions, Cauchy--Schwarz gives
\(\Phi_S(z+b)=\Phi_S(z)+Sb\) wherever the selected row limits exist,
for every \(b\in\ell^2\).
The affine orbit has second-moment bound
\(\|T\|_{\mathrm{op}}^2+\|b\|_2^2\).
This translation identity and the measurable composition lemma show
that the inverse affine map pushes the orbit back to \(\mu\).
Its pushforward of \(\mu\) is consequently also equivalent to \(\mu\).
Equivalence and the lower support bound show, by conditioning
finite row sums as before, that both linear parts have nonnegative
entries: a negative coefficient would produce an unbounded lower
tail, which no finite shift can repair.
Since \(T\) and \(T^{-1}\) have nonnegative matrix entries, the
off-diagonal inverse identities used above show that \(T\) is a
weighted permutation operator.
Equality of the scalar lower endpoints forces
\(-\sqrt\alpha c_i+b_i=-\sqrt\alpha\).
The product criterion \eqref{eq:restricted-gamma-scale-product}
then proves necessity and sufficiency.
\end{proof}

\subsection{Boundary behavior of the density and summability criteria}
\label{sec:restricted-threshold}

Let \(\rho\) be uniform on \([-\sqrt3,\sqrt3]\), let
\(\mu_2=\rho^{\otimes2}\), and consider \(T=\bigoplus_kO_{\theta_k}\).
Write
\(\delta_k=\min_{m\in\mathbb Z}|\theta_k-m\pi/2|\in[0,\pi/4]\).
For one block the affinity is the relative area of intersection
of a square and its rotation.
Removing the four corner triangles gives, for \(0<\delta\leq\pi/4\),
\begin{equation}
 \begin{aligned}
 \Aff(\mu_2,(O_\delta)_\#\mu_2)
 &=1-\frac{(\cos\delta+\sin\delta-1)^2}
                {2\sin\delta\cos\delta},\\
 1-\Aff(\mu_2,(O_\delta)_\#\mu_2)
 &=\frac{\delta}{2}+O(\delta^2).
 \end{aligned}
 \label{eq:restricted-uniform-rotation}
\end{equation}
At \(\delta=0\), the affinity is one.  Each nonzero \(\delta\) leaves positive area on both sides of the
support difference.
Consequently \cref{thm:block-kakutani} gives the complete trichotomy
\begin{equation}
 \begin{array}{c|l}
  \delta_k=0\text{ for every }k&\nu_T=\mu\\
  0<\sum_k\delta_k<\infty&
      \nu_T\not\perp\mu,\quad \nu_T\not\ll\mu,\quad\mu\not\ll\nu_T\\
  \sum_k\delta_k=\infty&\nu_T\perp\mu.
 \end{array}
 \label{eq:restricted-uniform-classification}
\end{equation}
In particular, \(\theta_k=1/k\) gives \(T-I\in\mathcal S_2\) but
\(\nu_T\perp\mu\).
With \(\theta_k=1/k^2\), the laws instead have positive affinity
and neither is absolutely continuous with respect to the other.

For comparison, use the variance-one rescaling of a density proportional
to \((1-x^2)^r\mathbf1_{\{|x|<1\}}\), \(r>1\).
Its zero-extended square root satisfies the Sobolev hypotheses.
Let \(f\) be this density, \(\rho(dx)=f(x)\,dx\),
\(\mu_2=\rho^{\otimes2}\), and \(g=\sqrt f\).
Writing \(J=4\|g'\|_2^2>1\), the skew-generator formula gives
\begin{equation}
 1-\Aff(\mu_2,(O_\theta)_\#\mu_2)=\frac{J-1}{4}\theta^2+o(\theta^2).
 \label{eq:restricted-smooth-rotation}
\end{equation}
Indeed, the planar skew generator has Frobenius norm squared two,
and \(h^2=\theta^2\langle A,\mathcal I A\rangle_F/4+o(\theta^2)\).
The only rotations preserving the square support are its signed
permutation symmetries.
Every rotated square shares a neighborhood of the origin with the
original square, on which both densities are positive; hence
\(\Aff(\mu_2,(O_\theta)_\#\mu_2)>0\) for every angle.
Affinity is continuous and strictly below one away from them,
so the three cases in \eqref{eq:restricted-uniform-classification}
now use \(\sum_k\delta_k^2\) in place of \(\sum_k\delta_k\).
Equivalence still requires every \(\delta_k=0\).
Thus the quadratic metric supplied by the Fourier/Fisher method
controls accumulation of Hellinger loss, while boundary compatibility
separately controls absolute continuity.

\subsection{An exact replacement for projected-energy equivalence}
\label{sec:restricted-projections}

For arbitrary zero sets of an unbounded density, the preceding support
theorems need not apply.
There is nevertheless a complete criterion in terms of finite marginals
which requires no dichotomy.
For densities satisfying \eqref{eq:restricted-sobolev},
\cref{thm:linear-equivalence-group} gives the corresponding operator
classification.  The following measure-theoretic criterion also applies
without Sobolev assumptions.

\begin{proposition}[Two-sided finite-marginal likelihood ratio tails]
\label{prop:restricted-projected-criterion}
Let \(\mu,\nu\) be any probability measures on \(\R^\N\), and let
\(\mu_n,\nu_n\) be their first \(n\) marginals.
Choose densities \(p_n,q_n\) with respect to any common dominating
measure \(m_n\).
For \(M>1\), define
\[
 D_n(M)=
 \int_{\{q_n>Mp_n\}}q_n\,dm_n+
 \int_{\{p_n>Mq_n\}}p_n\,dm_n.
\]
Then
\begin{equation}
 \mu\sim\nu
 \quad\Longleftrightarrow\quad
 \lim_{M\to\infty}\sup_nD_n(M)=0.
 \label{eq:restricted-ui-criterion}
\end{equation}
Equivalently, every finite pair is equivalent and its forward and
reverse likelihood-ratio martingales are uniformly integrable in their
respective reference measures.
Without any finite-marginal equivalence assumption,
\begin{equation}
 \Aff(\mu_n,\nu_n)\downarrow\Aff(\mu,\nu),\qquad
 \sup_n[-2\log\Aff(\mu_n,\nu_n)]<\infty
     \quad\Longleftrightarrow\quad \mu\not\perp\nu.
 \label{eq:restricted-energy-only-overlap}
\end{equation}
\end{proposition}
\begin{proof}
If the limit in \eqref{eq:restricted-ui-criterion} is zero, a set
where \(p_n=0<q_n\) must have zero \(\nu_n\)-mass for every \(n\);
the reverse statement follows from the second integral.
Hence \(\mu_n\sim\nu_n\).
Writing \(\ell_n(x)=q_n(\pi_n x)/p_n(\pi_n x)\), the two terms in \(D_n(M)\) are exactly
\[
 \E_\mu[\ell_n\mathbf1_{\{\ell_n>M\}}],
 \qquad
 \E_\nu[\ell_n^{-1}\mathbf1_{\{\ell_n^{-1}>M\}}].
\]
They are the uniform-integrability tests for the two mean-one
likelihood-ratio martingales.
Each \(L^1\) limit has mass one, and cylinder tests identify the
corresponding absolutely continuous measure.
This gives both directions of absolute continuity.
Conversely, under equivalence these martingales are the conditional
expectations of the full likelihood ratio and its reciprocal, so both
are uniformly integrable.

For the affinity assertion put \(m=(\mu+\nu)/2\),
\(u=d\mu/dm\), and \(v=d\nu/dm\).
Both densities are bounded by two.
For \(\mathcal F_n=\sigma(x_1,\ldots,x_n)\), set
\(u_n=\E_m[u\mid\mathcal F_n]\) and
\(v_n=\E_m[v\mid\mathcal F_n]\).
Martingale convergence and bounded convergence imply
\[
 \Aff(\mu_n,\nu_n)
       =\int\sqrt{u_nv_n}\,dm
       \longrightarrow\int\sqrt{uv}\,dm.
\]
Conditional Cauchy--Schwarz gives monotonicity.
Global affinity is zero exactly for mutually singular measures,
which proves \eqref{eq:restricted-energy-only-overlap}.
\end{proof}

For a density \(f\) and \(T\in\GL(H)\), each finite output
marginal still has a Lebesgue density: condition on the complement
of an invertible input minor as in
\cref{lem:finite-marginal-likelihood-ratios}.
That density may now vanish on a different set from
\(\prod_{i\leq n}f(x_i)\).
Thus \eqref{eq:restricted-ui-criterion} is an exact criterion for
arbitrary mixing even in this case.
Under \eqref{eq:restricted-sobolev},
\cref{thm:linear-equivalence-group} identifies the linear equivalence
group by a Hilbert--Schmidt condition and scalar row conditions for the
operator and its inverse.  This includes unbounded densities with gaps.
For bounded and one-sided laws,
\cref{thm:restricted-support-rigidity} further reduces the group to
explicit classes of coordinate permutations, without Sobolev assumptions.
